\section{Introduction}
\label{sec:intro}

Long-context language models are increasingly deployed in retrieval-augmented generation (RAG) and multi-document question answering, where models must identify and extract relevant information from extensive contexts. However, these models exhibit systematic position bias: they tend to over-attend to information at the beginning or end of the context while under-utilizing middle positions, a phenomenon known as ``lost-in-the-middle''~\citep{Liu2023LostIT}. This bias significantly degrades performance when relevant information appears in unfavorable positions.

Attention Sorting~\citep{Peysakhovich2023AttentionSC} addresses this limitation through iterative document reordering based on attention patterns. By measuring how much attention each document receives during generation and placing high-attention documents at the end (where recency-biased models attend more effectively), the method substantially improves long-context QA accuracy. However, Attention Sorting requires $k$ iterations of sorting, resulting in $k+1$ prefill passes per query. For $k=5$, this represents a 6$\times$ increase in prefill computation compared to vanilla generation, significantly increasing latency for long contexts.

We hypothesize that iterative sorting is primarily needed because raw attention scores are confounded by position bias: a relevant document far from the end may not receive sufficient attention in a single pass to be placed optimally. If we can explicitly estimate and subtract this position-dependent bias from attention scores, a single sorting pass should suffice to match iterative sorting accuracy. To test this hypothesis, we propose \textbf{Debiased One-Pass Attention Sorting}, which estimates a per-prompt position-bias curve from distractor documents and subtracts it to produce debiased relevance scores, requiring only 2 prefill passes.

Our experiments on two models with different bias characteristics refute this hypothesis. On LLaMA-2-7B-32K-Instruct, debiasing produces identical results to uncalibrated $k=1$ sorting (94.83\% accuracy), with zero improvement across 600 examples. On YaRN-Llama-2-7b-64k, which exhibits severe recency bias, debiasing improves accuracy by 8.67 percentage points over uncalibrated $k=1$ but remains 14.84pp behind $k=5$ sorting, closing only 37\% of the gap.

Our contributions are threefold. First, we propose and evaluate debiased one-pass attention sorting, a method that estimates and removes position bias from attention scores to enable single-pass document reordering. Second, we demonstrate that position bias is not the primary bottleneck limiting single-pass sorting on well-tuned models, as evidenced by zero improvement from debiasing on LLaMA-2-7B-32K-Instruct. Third, we identify that iterative sorting provides benefits beyond bias correction, likely from attention context refinement and error reduction, suggesting future work should investigate these mechanisms.
